\chapter*{Введение}
\addcontentsline{toc}{chapter}{Введение}
\label{ch:intro}

\textbf{Ansible} — система автоматизации управления конфигурациями и оркестрации,
работающая по принципу агентless: не требует установки агентов на управляемые узлы,
взаимодействуя с ними исключительно через протокол SSH.

Цель работы — развернуть Ansible на сервере \texttt{gateway}, настроить
безагентное подключение к двум клиентским машинам и запустить playbook,
собирающий данные мониторинга: IP-адрес, имя хоста, версию ОС и объём
свободного дискового пространства.

Задачи:
\begin{enumerate}
  \item Установить Ansible и его зависимости на \texttt{gateway}.
  \item Подготовить клиентские ВМ (\texttt{desktop1}, \texttt{wordpress}):
        установить \texttt{openssh-server} и \texttt{python3}.
  \item Сгенерировать SSH-ключ \texttt{ed25519} и распространить его на клиентов.
  \item Создать инвентарный файл \texttt{/etc/ansible/hosts} и конфигурацию
        \texttt{ansible.cfg}.
  \item Написать и запустить playbook \texttt{monitoring.yml},
        сохраняющий данные в \texttt{/etc/ansible/monitoring/}.
  \item Проверить результаты в файлах \texttt{*\_info.txt}.
\end{enumerate}

\textbf{Стенд:}
\begin{itemize}
  \item ВМ \texttt{gateway} — Ubuntu~20.04.5~LTS, два адаптера:
        \texttt{enp0s3}~(NAT) и \texttt{enp0s8}~(Internal,
        \texttt{192.168.\cfgLabStudentN.1/24});
        роль: Ansible control node.
  \item ВМ \texttt{desktop1} — Ubuntu~20.04.5~Desktop, адаптер Internal,
        \texttt{192.168.\cfgLabStudentN.10}; роль: Ansible client~1.
  \item ВМ \texttt{wordpress} — Ubuntu~20.04.5~Server, адаптер Internal,
        \texttt{192.168.\cfgLabStudentN.6}; роль: Ansible client~2.
\end{itemize}

\endinput
